\chapter{اللقاء التاسع: تابع تفسير سورة البقرة (٢)}

\section{مقدمة: كيف يواجه المسلم الابتلاءات والمحن؟}
السلام عليكم ورحمة الله وبركاته. الحمد لله رب العالمين، والصلاة والسلام على نبينا محمد وعلى آله وصحبه وسلم. 
أحببت أن أبدأ هذا اللقاء بسؤال يأتيني كثيراً من الشباب: كيف يمكن للمسلم أن يواظب على جدول أعماله اليومية النافعة وهو يمر بمشكلات في نفسه أو ولده، أو يرى ما ينزل بالمسلمين من الظلم والحوادث؟

أقول باختصار: أما التأثر فكل مسلم ينبغي أن يتأثر ويحزن لما يصيب المؤمنين من الشر، لكن أن يقعده هذا الحزن عن العمل أو يبعث في نفسه اليأس والقنوط، فهذا ليس صحيحاً؛ بل ينبغي أن يزيده قوة وعزماً. فإذا لم يعمل الإنسان في أوقات الشدائد فمتى يعمل؟ 

خير ما تواجه به المصيبة هو: \textbf{الصبر، والاحتساب، والعمل الصالح}. 
وتذكر أن هذا الابتلاء تستخرج به معاني الإيمان وترفع به الدرجات. ولنا في رسول الله صلى الله عليه وسلم وأصحابه أسوة حسنة؛ فلما توفي النبي صلى الله عليه وسلم وعظم الخطب، خرج \rawi{أبو بكر الصديق} رضي الله عنه وتلا قوله تعالى: \begin{ayah}وَمَا مُحَمَّدٌ إِلَّا رَسُولٌ قَدْ خَلَتْ مِن قَبْلِهِ الرُّسُلُ ۚ أَفَإِن مَّاتَ أَوْ قُتِلَ انقَلَبْتُمْ عَلَىٰ أَعْقَابِكُمْ ۚ وَمَن يَنقَلِبْ عَلَىٰ عَقِبَيْهِ فَلَن يَضُرَّ اللَّهَ شَيْئًا ۗ وَسَيَجْزِي اللَّهُ الشَّاكِرِينَ\end{ayah}. فكأن الناس لم يعلموا أن الله أنزل هذه الآية حتى تلاها. فالقرآن شفاء لما في الصدور، وكلما نزل بك بلاء فاقرأ الآيات وتذكر ما فيها من المعاني والهدى.

وهذه المقدمة الوعظية ليست غريبة عن جو الدرس؛ لأن الشيخ أراد بها أن يربط طالب العلم من أول الأمر بأن الاشتغال بالوحي ليس عملاً ذهنياً مجرداً، بل هو من أعظم ما يثبت القلب وقت الفتن والشدائد.

\separator

\section{تأويل (يا أيها الناس اعبدوا ربكم)}
وصلنا بحمد الله إلى قوله تعالى: \begin{ayah}يَا أَيُّهَا النَّاسُ اعْبُدُوا رَبَّكُمُ الَّذِي خَلَقَكُمْ وَالَّذِينَ مِن قَبْلِكُمْ لَعَلَّكُمْ تَتَّقُونَ\end{ayah}.

\isnad{قال الإمام الطبري رحمه الله:} «أمر جل ثناؤه الفريقين اللذين أخبر عن أحدهما أنه سواء عليهم أأنذرتهم أم لم تنذرهم... وعن الآخر أنه يخادع الله والذين آمنوا... وغيرهم من سائر خلقه المكلفين، بالاستكانة والخضوع له بالطاعة، وإفراد الربوبية له والعبادة دون الأوثان...».

\begin{fawaid}[ضابط الخطاب المكي والمدني]
هذه الآية فيها تخطئة لمن زعم من المتأخرين أن كل سورة فيها خطاب بـ (يا أيها الناس) فهي مكية، وأن (يا أيها الذين آمنوا) لا تكون إلا في المدنية. فسورة البقرة مدنية بالإجماع، وفيها نداء بـ (يا أيها الناس).
\end{fawaid}

\subsection{مسألة تكليف ما لا يطاق}
\isnad{قال الطبري في هذه الآية:} «وهذه الآية من أدل الدليل على فساد قول من زعم أن تكليف ما لا يطاق إلا بمعونة الله غير جائز... وذلك أن الله جل وعز أمر من وصفنا بعبادته والتوبة من الكفر بعد إخباره عنهم أنهم لا يؤمنون».

وقد سبق أن بينا أن هذا الموضع محل نقد؛ فالطبري يظن أن إخبار الله عن الكفار بأنهم لن يؤمنوا (لسابق علمه بهم) يُعد من باب تكليفهم بما لا يطيقون. والصواب عند أهل السنة أن الله لا يكلف نفساً إلا وسعها (من جهة القدرة والأهلية)، أما كون العبد لن يفعل المأمور به في علم الله فهذا لا يُسمى "تكليفاً بما لا يطاق" لاستطاعة العبد له ابتداءً.

\separator

\section{تأويل (لعلكم تتقون) ونقد منهج اللغويين في الإشكالات}
\isnad{قال الطبري:} «فإن قال لنا قائل: وكيف قال جل ثناؤه (لعلكم تتقون) أو لم يكن عالما بما يصير إليه أمرهم... حتى أخرج الخبر مخرج الشك؟ قيل: ذلك على غير المعنى الذي توهمت، وإنما معنى ذلك: اعبدوا ربكم لتتقوه بطاعته».

\begin{fawaid}[خطر جعل القواعد النحوية الضيقة حاكمة على القرآن]
نلاحظ هنا منهج \rawi{الطبري} في طرح الإشكالات التي أثارها بعض النحويين (كالفراء والأخفش وأبي عبيدة). فقد يضيق أحدهم معنى أداة من الأدوات، ثم إذا وجد القرآن أوسع من تقريره توهم الإشكال. والصواب أن القرآن نفسه من أعلى شواهد العربية، فلا يصح أن نجعل قاعدة مستنبطة من بعض الاستعمالات حاكمة عليه.
\end{fawaid}

\separator

\section{تأويل (فلا تجعلوا لله أنداداً وأنتم تعلمون)}
\isnad{قال الطبري:} «والأنداد جمع نِد، والند: العِدْل والمثل... ونهاهم الله أن يشركوا به شيئاً».

ثم تطرق لتأويل \begin{ayah}وَأَنتُمْ تَعْلَمُونَ\end{ayah}. واختلف المفسرون فيمن عُني بذلك؛ هل هم مشركو العرب وأهل الكتاب معاً، أم أهل الكتاب خاصة؟

\isnad{قال الطبري:} «وأحسب أن الذي دعا \rawi{مجاهداً} إلى إضافته لأهل الكتاب دون غيرهم، الظن منه بالعرب أنها لم تكن تعلم أن الله خالقها ورازقها... ولكن الله أخبر عنها أنها كانت تقر بوحدانيته... ولئن سألتهم من خلقهم ليقولن الله».

\begin{fawaid}[التثبت من مأخذ القول قبل نقده]
الذي يظهر أن \rawi{الطبري} افترض أن \rawi{مجاهداً} خص الآية بأهل الكتاب لأنه يظن أن العرب كانت تجهل الخالق. وهذا افتراض غير دقيق؛ فمجاهد أجلّ من أن يجهل إقرار العرب بربوبية الله (وهو الذي قرأ القرآن على ابن عباس وسأله عن كل آية). ومجاهد لم يصرح بحصر الآية في أهل الكتاب، بل ذكرهم كمثال لدخولهم في المعنى.
\end{fawaid}

\separator

\section{تأويل (وإن كنتم في ريب مما نزلنا على عبدنا فأتوا بسورة من مثله)}
\isnad{قال أبو جعفر:} «وهذا من الله جل ثناؤه احتجاج لنبيه محمد صلى الله عليه وسلم على مشركي قومه ومنافقيهم... فإن كنتم في شك مما نزلنا على عبدنا من النور والبرهان فأتوا بحجة تدفع حجته».

ثم اختلفوا في الضمير في قوله: (من مثله)؛ هل يعود على القرآن أم يعود على محمد صلى الله عليه وسلم (بمعنى: فأتوا بسورة من رجل أمي مثل محمد)؟
ورجح \rawi{الطبري} أن الضمير يعود على القرآن، واستدل بآية أخرى لرد القول المخالف، وهي قوله تعالى: \begin{ayah}أَمْ يَقُولُونَ افْتَرَاهُ ۖ قُلْ فَأْتُوا بِسُورَةٍ مِّثْلِهِ\end{ayah}. فهنا صرحت الآية بـ (مثله) أي مثل القرآن دون حرف (مِن). وهذا من أجمل الترجيحات بتفسير القرآن بالقرآن.

\subsection{تأويل (وادعوا شهداءكم)}
رجح \rawi{الطبري} أن المراد بالشهداء هنا (الأعوان والأنصار)، ورَدَّ قول \rawi{ابن جريج} الذي فسرها بـ (الذين يشهدون لكم أنه مثل القرآن). وحجة الطبري في رده: أن الناس حينها إما مؤمن صحيح أو كافر أو منافق، فمن الذي سيشهد للكفار بالحق إن كانوا يبطلون الحق؟ وهذا استدلال عقلي قوي من الطبري في الترجيح.

\begin{fawaid}[الترجيح بتقسيم الواقع والمخاطبين]
من دقائق منهج \rawi{الطبري} هنا أنه لا يكتفي بمجرد الاحتمال اللغوي، بل ينظر في واقع المخاطبين وأحوالهم: من الذي يمكن أن يدخل تحت هذا اللفظ في ذلك السياق؟ وهذا باب قوي في إبطال بعض الأقوال التي تصح لغة ولا تستقيم واقعاً وسياقاً.
\end{fawaid}

\separator

\section{تأويل (فاتقوا النار التي وقودها الناس والحجارة)}
\isnad{قال أبو جعفر:} «فأخبرهم أن الناس وقودها وأن الحجارة وقودها... يعني حطبها».
ثم ساق آثاراً عن \rawi{ابن مسعود} و\rawi{ابن عباس} و\rawi{مجاهد} أن الحجارة هنا هي "حجارة من كبريت أسود".

\begin{fawaid}[الاكتفاء بالقرآن في الموعظة وترك الروايات الضعيفة]
هذه الآثار المروية في تفصيل نوع حجارة جهنم هي من أخبار الغيب التي لا يعتمد عليها إذا لم تثبت عن النبي صلى الله عليه وسلم (وأغلبها مأخوذ من أهل الكتاب). والداعية والمربي لا يحتاج إلى هذه التفاصيل الواهية أو الغريبة لتخويف الناس من النار؛ ففي آيات القرآن البينة والأحاديث الصحيحة الغنى والكفاية للترغيب والترهيب دون تكلف.
\end{fawaid}

\separator

\section{تأويل (كلما رزقوا منها من ثمرة رزقاً قالوا هذا الذي رزقنا من قبل)}
اختلفوا في معنى: (هذا الذي رزقنا من قبل). 
فرأى بعضهم كـ \rawi{يحيى بن أبي كثير} أنها تعني: رزقنا من قبل في \textbf{الجنة} (لتشابه الثمار هناك)، ورأى آخرون كـ \rawi{ابن عباس} و\rawi{مجاهد} أنها تعني: رزقنا من قبل في \textbf{الدنيا}.

ورجح \rawi{الطبري} بقوة القول الثاني (أنهم قصدوا ما رزقوه في الدنيا)، محتجاً بأن الله عمم ذلك بقوله (كلما رزقوا)، فلو كان المعنى أنهم يقصدون ثماراً أكلوها في الجنة سابقاً، فكيف سيقولون ذلك عند \textbf{أول} ثمرة يأكلونها حين دخولهم الجنة؟! فبطل القول الأول.

\begin{fawaid}[رد تعنت النحويين في الاستعمال المجازي]
تذرع بعض النحويين بأنهم لا يمكن أن يقصدوا (ما رزقوه في الدنيا) لأن تلك الثمار فنيت وانعدمت. فأجاب الطبري بأن هذا التعنت خلاف مخرج كلام العرب؛ فالعربي يقول لمن أعد له طعاماً كطعامه: "هذا طعامي في منزلي"، ولا يتوهم سامع أنه يقصد عين الطعام، بل نوعه وجنسه.
\end{fawaid}

\begin{fawaid}[العموم في اللفظ قد يهدم بعض الأقوال من أصلها]
استدلال \rawi{الطبري} بقوله تعالى: \begin{ayah}كُلَّمَا رُزِقُوا\end{ayah} من أنفع أمثلة الاحتجاج بعموم اللفظ؛ إذ بيّن أن التعميم يشمل أول رزق في الجنة، فيبطل بذلك تفسير من قال إنهم يقصدون ثمرة سابقة في الجنة نفسها.
\end{fawaid}

\separator

\section{تأويل (إن الله لا يستحيي أن يضرب مثلاً ما بعوضة فما فوقها)}
ذكر \rawi{الطبري} الخلاف في سبب نزول هذه الآية؛ هل هي للرد على المنافقين لإنكارهم الأمثال السابقة في نفس السورة (مَثَلُهُمْ كَمَثَلِ الَّذِي اسْتَوْقَدَ نَاراً)، أم هي رداً على المشركين الذين أنكروا ذكر الله للعنكبوت والذباب في السور الأخرى؟

ورجح \rawi{الطبري} القول الأول، لقوة ارتباطه بسياق السورة واتصال الآيات. وهذه من قواعد الطبري الجليلة: (تقديم ما يعززه سياق السورة على ما يبعد عنه).

\separator

\section{تأويل (وما يضل به إلا الفاسقين)}
\isnad{قال أبو جعفر:} «يعني وما يضل الله بالمثل الذي يضربه إلا التاركين اتباع أمره من أهل النفاق».

\begin{fawaid}[أسباب الهداية والإضلال]
هذه الآية من أهم الآيات في تقرير باب "القدر" للرد على الجبرية؛ فهي تثبت أن الله تبارك وتعالى لا يُضل عبداً إلا بسببٍ من العبد وكسبه، فالإضلال كان نتيجة لفسقهم المسبق، كما قال تعالى: \begin{ayah}فَلَمَّا زَاغُوا أَزَاغَ اللَّهُ قُلُوبَهُمْ\end{ayah}.
\end{fawaid}

\separator

\section{تأويل (كَيْفَ تَكْفُرُونَ بِاللَّهِ وَكُنتُمْ أَمْوَاتًا فَأَحْيَاكُمْ)}
ساق \rawi{الطبري} الآثار عن \rawi{ابن مسعود} و\rawi{ابن عباس} في تفسيرها بأن الموتة الأولى هي كونهم عدماً قبل أن يُخلقوا، ثم أحياهم، ثم يميتهم الموتة المعهودة، ثم يحييهم للبعث.
وهذا من أروع أمثلة (تفسير القرآن بالقرآن)، فقد فسرها المفسرون بقوله تعالى في سورة غافر: \begin{ayah}قَالُوا رَبَّنَا أَمَتَّنَا اثْنَتَيْنِ وَأَحْيَيْتَنَا اثْنَتَيْنِ\end{ayah}.

\separator

\section{تأويل (ثُمَّ اسْتَوَى إِلَى السَّمَاءِ) ونقض تأويلات المعطلة}
\isnad{قال أبو جعفر:} «والاستواء في كلام العرب منصرف على وجوه... وأولى المعاني بقول الله (ثم استوى إلى السماء فسواهن) علا عليهن وارتفع فدبرهن بقدرته...».

ثم شن \rawi{الطبري} هجوماً قوياً على المعتزلة والمعطلة الذين ينكرون صفة "العلو" ويفرون إلى تأويلها، \isnad{فقال:} «والعجب ممن أنكر المعنى المفهوم من كلام العرب في تأويل قول الله (ثم استوى إلى السماء) الذي هو بمعنى العلو والارتفاع، هرباً عند نفسه من أن يلزمه بزعمه... أن يكون إنما علا وارتفع بعد أن كان تحتها... ثم لم ينجُ مما هرب منه؛ فيقال له: أزعمت أن تأويل استوى: أقبل؟ أفكان مدبراً عن السماء فأقبل إليها؟!».

\begin{fawaid}[إثبات صفات الأفعال ومنهج الإلزام]
من مقاصد الطبري الكبرى في هذا الكتاب إبطال منهج المعتزلة في نفي الصفات. وقد استخدم معهم قوة "الإلزام" العقلي؛ فبين لهم أن ما يفرون منه من لوازم باطلة في إثبات (العلو)، سيلزمهم أشنع منه في اللفظ الذي أولوا إليه (أقبل)! فالأسلم والأحق إثبات ما أثبته الله لنفسه على وجه الكمال دون تكييف أو تعطيل.
\end{fawaid}

\separator

\section{تأويل (إِنِّي جَاعِلٌ فِي الْأَرْضِ خَلِيفَةً)}
اختلف المفسرون في معنى الخليفة هنا؛ فقيل: يخلفون من كان قبلهم من الجن الذين سكنوا الأرض، وقيل: بل هم ولد آدم الذين يخلف بعضهم بعضاً قرناً بعد قرن.

\begin{fawaid}[نقد مقولة: آدم خليفة الله في الأرض]
نُسب لبعضهم (كابن زيد) أن آدم خليفة (عن الله) في الأرض. وقد اشتد إنكار المحققين كشيخ الإسلام \rawi{ابن تيمية} على هذا الإطلاق؛ فالله تعالى هو الخالق المدبر الشهيد الذي لا يغيب، والمخلوق إنما يستخلف غيره لعجزه أو غيابه. فآدم وذريته يخلف بعضهم بعضاً، وليسوا خلفاء "عن الله" عز وجل.
\end{fawaid}

\subsection{تأويل قول الملائكة: (أَتَجْعَلُ فِيهَا مَن يُفْسِدُ فِيهَا)}
تساءل الطبري: كيف علمت الملائكة أنهم سيفسدون في الأرض ولم يُخلقوا بعد؟ هل هو رجم بالغيب؟
أورد الطبري أثراً طويلاً عن \rawi{السدي} يذكر فيه تفاصيل (إسرائيلية) لخلق آدم وإبليس والملائكة، وخلاصته أن الله هو الذي أخبر الملائكة أن ذرية هذا الخليفة سيفسدون، فتعجبت الملائكة وسألت: (أتجعل فيها من يفسد فيها).

وقد انتقد الطبري بعض الروايات التي تناقض ظواهر القرآن، مقرراً قاعدته العظيمة: \textbf{«غير جائز أن يترك المفهوم من ظاهر الكتاب... إلى باطل لا دلالة عليه من ظاهر التنزيل ولا خبر عن الرسول صلى الله عليه وسلم ولا فيه من الحجة إجماع مستفيض»}.

\separator

بهذا نكون قد أتممنا بحمد الله هذا القدر من تفسير سورة البقرة، وسنكمل في اللقاء القادم بإذن الله. والسلام عليكم ورحمة الله وبركاته.
